\section{משפט שטורם-ליוביל, פונקציית גרין ומשוואות דיפרנציאליות אי-הומוגניות — 18 בנובמבר 2025}

\subsection*{סקירה: משפט שטורם-ליוביל}

בשיעור הקודם נחשפנו למשפט שטורם-ליוביל, אחד מאבני היסוד החשובות ביותר בתורת המשוואות הדיפרנציאליות הרגילות. משפט זה עוסק במשוואות דיפרנציאליות ממעלה שנייה מהצורה
\[
L u(x) = \lambda u(x),
\]
כאשר \(L\) הוא אופרטור דיפרנציאלי לינארי ממעלה שנייה המוגדר על ידי
\[
L u = \frac{d}{dx}\left[p_0(x) \frac{du}{dx}\right] + p_1(x) \frac{du}{dx} + p_2(x) u(x).
\]
האופרטור \(L\) נדרש להיות \textbf{צמוד לעצמו} (self-adjoint), כלומר מקיים את התנאי \(p_0'(x) = p_1(x)\). כאשר מתלווים לבעיה זו תנאי שפה הומוגניים בקצות הקטע \([a,b]\) (תנאי דיריכלה, נוימן או מעורבים), משפט שטורם-ליוביל מבטיח לנו שלוש תוצאות יסודיות:

\paragraph{תכונה 1: ערכים עצמיים ממשיים.}
כל הערכים העצמיים \(\lambda_i\) הם מספרים ממשיים. מבחינה פיזיקלית, זוהי תכונה חיונית: ערכים עצמיים מייצגים תדרי תנודה, רמות אנרגיה, או קבועי התפשטות, וכולם חייבים להיות ממשיים במערכות פיזיקליות ריאליות.

\paragraph{תכונה 2: אורתוגונליות של הפונקציות העצמיות.}
הפונקציות העצמיות \(u_i(x)\) ו-\(u_j(x)\) אורתוגונליות זו לזו במכפלה הפנימית המשוקללת:
\[
\langle u_i, u_j \rangle = \int_a^b u_i(x) u_j(x) w(x) \, dx = 0 \quad \text{אם } i \neq j,
\]
כאשר \(w(x)\) היא פונקציית המשקל. במקרה שבו \(i = j\), האינטגרל נותן את הנורמה (או 1 במקרה מנורמל). תכונה זו מאפשרת לנו לבנות בסיס אורתונורמלי במרחב הפונקציות.

\paragraph{תכונה 3: פריסה של כל פונקציה בטור פונקציות עצמיות.}
כל פונקציה \(f(x)\) בקטע \([a,b]\) ניתנת לפריסה יחידה בטור של הפונקציות העצמיות:
\[
f(x) = \sum_{n=1}^\infty c_n u_n(x).
\]
המקדמים \(c_n\) נקבעים באופן יחיד על ידי:
\[
c_n = \frac{\langle f, u_n \rangle}{\langle u_n, u_n \rangle} = \frac{\int_a^b f(x) u_n(x) w(x) \, dx}{\int_a^b u_n^2(x) w(x) \, dx}.
\]

משפט זה חזק במיוחד משום שהוא מבטיח שכאשר נפתור משוואות דיפרנציאליות חלקיות (כמו משוואת החום, משוואת הגלים, או משוואת לפלס), תמיד נגיע לתת-בעיה מסוג שטורם-ליוביל. הפריסה היחידה מבטיחה שהפתרון שנמצא הוא הפתרון היחיד, מה שמקנה כוח עצום לשיטה זו.

\subsection*{הבאת אופרטור להיות צמוד לעצמו}

לא כל אופרטור דיפרנציאלי צמוד לעצמו באופן טבעי. נניח שניתנה לנו משוואה:
\[
L u + \lambda w_1(x) u = 0,
\]
כאשר
\[
L u = p_0(x) u''(x) + p_1(x) u'(x) + p_2(x) u(x),
\]
ו-\(L\) \textbf{אינו} צמוד לעצמו, כלומר \(p_0'(x) \neq p_1(x)\). המטרה שלנו היא להביא את האופרטור למצב שבו הוא כן צמוד לעצמו, כך שנוכל להשתמש במשפט שטורם-ליוביל.

הרעיון הוא להכפיל את כל המשוואה בפונקציה \(F(x)\) מתאימה:
\[
F(x) \left[ p_0(x) u''(x) + p_1(x) u'(x) + p_2(x) u(x) \right] + \lambda F(x) w_1(x) u(x) = 0.
\]
כעת נדרוש שהמשוואה החדשה תהיה צמודה לעצמה, כלומר:
\[
(F p_0)' = F p_1.
\]
פיתוח של הדרישה הזו נותן:
\[
F' p_0 + F p_0' = F p_1,
\]
ולכן
\[
F' p_0 = F(p_1 - p_0').
\]
זו משוואה דיפרנציאלית עבור \(F\):
\[
\frac{F'}{F} = \frac{p_1 - p_0'}{p_0}.
\]
פתרון משוואה זו (אינטגרציה של שני האגפים) נותן:
\[
\ln F = \int_a^x \frac{p_1(x') - p_0'(x')}{p_0(x')} \, dx',
\]
ולכן
\[
F(x) = \exp\left( \int_a^x \frac{p_1(x') - p_0'(x')}{p_0(x')} \, dx' \right).
\]

שימו לב: גבול האינטגרציה התחתון (\(a\)) הוא שרירותי, משום שקבוע כפל ב-\(F\) לא משפיע על המשוואה ההומוגנית. פונקציית המשקל החדשה היא \(w(x) = F(x) w_1(x)\), והיא חיובית בכל הקטע (כי \(F\) מעריכית חיובית, ו-\(w_1\) חיובית לפי ההנחה המקורית).

\paragraph{בדיקה:} אם האופרטור היה כבר צמוד לעצמו (כלומר \(p_0' = p_1\)), אז \(F(x) = 1\), ולא השתנה דבר.

\subsection*{דוגמה: משוואת לגר}

משוואת לגר (Laguerre) מופיעה במכניקה קוונטית כאשר פותרים את פונקציות הגל של אטום המימן. בצורתה הפשוטה היא נראית כך:
\[
x u''(x) + (1 - x) u'(x) + \lambda u(x) = 0, \quad x \in [0, \infty).
\]
כאן \(p_0(x) = x\), \(p_1(x) = 1 - x\), ו-\(p_2(x) = 0\). בדיקה מהירה מראה ש-\(p_0'(x) = 1 \neq p_1(x)\), ולכן המשוואה אינה צמודה לעצמה. כדי לתקן זאת, נחשב:
\[
\frac{p_1 - p_0'}{p_0} = \frac{(1-x) - 1}{x} = \frac{-x}{x} = -1.
\]
לכן
\[
F(x) = \exp\left( \int_0^x (-1) \, dx' \right) = e^{-x}.
\]
נכפיל את המשוואה כולה ב-\(e^{-x}\):
\[
e^{-x} x u''(x) + e^{-x}(1-x) u'(x) + \lambda e^{-x} u(x) = 0.
\]
ניתן לרשום זאת בצורה:
\[
\frac{d}{dx}\left[ e^{-x} x u'(x) \right] + \lambda e^{-x} u(x) = 0.
\]
כעת האופרטור צמוד לעצמו, ופונקציית המשקל החדשה היא \(w(x) = e^{-x}\). תחת פונקציית משקל זו, הפולינומים של לגר אורתוגונליים.

\subsection*{מעבר למשוואות אי-הומוגניות: פונקציית גרין}

לאחר שהבנו כיצד לפתור בעיות הומוגניות, נרצה כעת לטפל במשוואות אי-הומוגניות מהצורה:
\[
L u(x) - \Lambda u(x) = f(x), \quad x \in [a,b],
\]
עם תנאי שפה הומוגניים. כאן \(L\) הוא אופרטור צמוד לעצמו, \(\Lambda\) פרמטר ממשי נתון (שונה מהערכים העצמיים \(\lambda_n\) של הבעיה ההומוגנית), ו-\(f(x)\) כוח מאלץ נתון.

\paragraph{זיכרו:} הבעיה ההומוגנית המקבילה היא
\[
L u_n(x) = \lambda_n u_n(x),
\]
ופתרנו אותה ומצאנו פונקציות עצמיות \(u_n(x)\) וערכים עצמיים \(\lambda_n\).

\paragraph{רעיון הפתרון:}
מכיוון שהפונקציות העצמיות \(u_n(x)\) יוצרות בסיס שלם, נוכל לפרוס גם את \(u(x)\) וגם את \(f(x)\) בטור:
\[
u(x) = \sum_{n=1}^\infty c_n u_n(x), \quad f(x) = \sum_{n=1}^\infty d_n u_n(x).
\]
המקדמים \(d_n\) ידועים (נתונים על ידי \(f\)), והמקדמים \(c_n\) לא ידועים. נציב את הפריסות במשוואה המקורית:
\[
L \left( \sum_{n=1}^\infty c_n u_n \right) - \Lambda \left( \sum_{n=1}^\infty c_n u_n \right) = \sum_{n=1}^\infty d_n u_n.
\]
מכיוון ש-\(L u_n = \lambda_n u_n\), מקבלים:
\[
\sum_{n=1}^\infty c_n \lambda_n u_n - \Lambda \sum_{n=1}^\infty c_n u_n = \sum_{n=1}^\infty d_n u_n,
\]
או
\[
\sum_{n=1}^\infty c_n (\lambda_n - \Lambda) u_n = \sum_{n=1}^\infty d_n u_n.
\]
מכיוון שהפונקציות העצמיות בלתי תלויות לינארית, כל מקדם חייב להתאפס בנפרד:
\[
c_n (\lambda_n - \Lambda) = d_n,
\]
ולכן
\[
c_n = \frac{d_n}{\lambda_n - \Lambda}.
\]

\paragraph{תנאי הכרחי:}
הפתרון קיים ויחיד רק אם \(\Lambda \neq \lambda_n\) לכל \(n\). אם \(\Lambda = \lambda_k\) עבור \(k\) כלשהו, ו-\(d_k \neq 0\), אין פתרון (תנאי רזוננס). אם \(d_k = 0\), יש אינסוף פתרונות (המקדם \(c_k\) שרירותי).

\subsection*{הגדרת פונקציית גרין — הצגה ראשונה}

לאחר שמצאנו את \(c_n\), הפתרון הוא:
\[
u(x) = \sum_{n=1}^\infty c_n u_n(x) = \sum_{n=1}^\infty \frac{d_n}{\lambda_n - \Lambda} u_n(x).
\]
נזכור ש-\(d_n\) ניתן על ידי:
\[
d_n = \int_a^b u_n^*(x') f(x') \, dx'.
\]
(בהנחה שהפונקציות מנורמלות; אחרת צריך לחלק בנורמה). נציב:
\[
u(x) = \sum_{n=1}^\infty \frac{u_n(x)}{\lambda_n - \Lambda} \int_a^b u_n^*(x') f(x') \, dx'.
\]
כעת נגדיר את \textbf{פונקציית גרין} \LR{(Green's function): }
\[
G(x, x'; \Lambda) = \sum_{n=1}^\infty \frac{u_n(x) u_n^*(x')}{\lambda_n - \Lambda}.
\]
בעזרת הגדרה זו, הפתרון מתבטא באופן אלגנטי:
\[
u(x) = \int_a^b G(x, x'; \Lambda) f(x') \, dx'.
\]

\paragraph{הכוח של הפורמליזם:}
פונקציית גרין \(G(x, x'; \Lambda)\) תלויה רק בבעיה ההומוגנית (בפונקציות העצמיות ובערכים העצמיים). ברגע שמצאנו אותה פעם אחת, כל כוח מאלץ \(f(x)\) חדש מוביל לפתרון מיידי על ידי חישוב האינטגרל הנ"ל. אין צורך לפתור מחדש את המשוואה הדיפרנציאלית.

\paragraph{תכונה: סימטריה.}
אם הפונקציות העצמיות ממשיות (\(u_n^* = u_n\)), אז
\[
G(x, x'; \Lambda) = G(x', x; \Lambda),
\]
כלומר פונקציית גרין סימטרית בהחלפת \(x \leftrightarrow x'\).

\subsection*{הצגה שנייה לפונקציית גרין: גישה ישירה}

ההצגה הראשונה של \(G\) כטור אינסופי יפה מבחינה פורמלית, אך קשה לחישוב כאשר לא ניתן למצוא את הפונקציות העצמיות באופן אנליטי. נפתח כעת הצגה שנייה, שמתאימה למקרים בהם אנחנו פותרים את המשוואה הדיפרנציאלית ישירות.

הרעיון: נפתור את המשוואה
\[
L u(x) - \Lambda u(x) = \delta(x - x'),
\]
כאשר \(\delta(x - x')\) היא פונקציית דלתא של דיראק. הפתרון של משוואה זו הוא בדיוק \(G(x, x')\), כפי שניתן לראות על ידי הצבה בביטוי האינטגרלי:
\[
u(x) = \int_a^b G(x, x'; \Lambda) \delta(x' - x_0) \, dx' = G(x, x_0; \Lambda).
\]

\paragraph{תכונות של \(G\):}
\begin{enumerate}
\item \textbf{רציפות:} \(G(x, x')\) רציפה ב-\(x = x'\). הסיבה: אם \(G\) הייתה בעלת אי-רציפות סופית, הנגזרת הראשונה שלה הייתה מכילה \(\delta\), והנגזרת השנייה \(\delta'\). אך במשוואה מופיעה רק \(\delta\), לא \(\delta'\), ולכן \(G\) חייבת להיות רציפה.

\item \textbf{קפיצה בנגזרת:} הנגזרת \(\frac{\partial G}{\partial x}\) קופצת ב-\(x = x'\). נחשב את הקפיצה על ידי אינטגרציה של המשוואה הדיפרנציאלית:
\[
\int_{x' - \epsilon}^{x' + \epsilon} \left[ \frac{d}{dx}\left( p_0(x) \frac{dG}{dx} \right) + q(x) G - \Lambda G \right] dx = \int_{x' - \epsilon}^{x' + \epsilon} \delta(x - x') \, dx = 1.
\]
כאשר \(\epsilon \to 0\), האינטגרלים של \(G\) ו-\(q G\) ו-\(\Lambda G\) שואפים לאפס (כי \(G\) רציפה), ונשאר:
\[
p_0(x') \left[ \frac{\partial G}{\partial x}\bigg|_{x = x'^+} - \frac{\partial G}{\partial x}\bigg|_{x = x'^-} \right] = 1,
\]
כלומר
\[
\left. \frac{\partial G}{\partial x} \right|_{x'^+} - \left. \frac{\partial G}{\partial x} \right|_{x'^-} = \frac{1}{p_0(x')}.
\]
\end{enumerate}

\subsection*{בניית פונקציית גרין: שיטת התפירה}

נפתור את המשוואה ההומוגנית בשני הקטעים \([a, x')\) ו-\((x', b]\) בנפרד:
\begin{align*}
L y_1(x) - \Lambda y_1(x) &= 0, \quad x \in [a, x'), \\
L y_2(x) - \Lambda y_2(x) &= 0, \quad x \in (x', b].
\end{align*}
לכל אחת מהמשוואות הללו יש פתרון כללי עם שני קבועים. נשתמש בתנאי השפה: \(y_1\) מקיימת את תנאי השפה ב-\(x = a\), ו-\(y_2\) מקיימת את תנאי השפה ב-\(x = b\). זה מותיר קבוע אחד לא ידוע בכל אחד מהפתרונות:
\[
G(x, x') = \begin{cases}
C_1 y_1(x), & x < x', \\
C_2 y_2(x), & x > x'.
\end{cases}
\]
הקבועים \(C_1\) ו-\(C_2\) נקבעים על ידי שני תנאי התפירה ב-\(x = x'\):
\begin{enumerate}
\item רציפות: \(C_1 y_1(x') = C_2 y_2(x')\),
\item קפיצה בנגזרת: \(C_2 y_2'(x') - C_1 y_1'(x') = \frac{1}{p_0(x')}\).
\end{enumerate}

זוהי מערכת של שתי משוואות לינאריות בשני נעלמים \(C_1, C_2\). נפתור אותה באמצעות כלל קרמר. נגדיר:
\[
\Delta = \begin{vmatrix}
y_1(x') & -y_2(x') \\
-y_1'(x') & y_2'(x')
\end{vmatrix} = y_1(x') y_2'(x') - y_1'(x') y_2(x') = W(x'),
\]
כאשר \(W(x')\) הוא הרונסקיאן (Wronskian) של \(y_1\) ו-\(y_2\).

הפתרונות הם:
\[
C_1 = \frac{1}{p_0(x') W(x')} \cdot \begin{vmatrix}
0 & -y_2(x') \\
\frac{1}{p_0(x')} & y_2'(x')
\end{vmatrix} = \frac{y_2(x')}{p_0(x') W(x')},
\]
\[
C_2 = \frac{1}{p_0(x') W(x')} \cdot \begin{vmatrix}
y_1(x') & 0 \\
-y_1'(x') & \frac{1}{p_0(x')}
\end{vmatrix} = \frac{y_1(x')}{p_0(x') W(x')}.
\]

לכן ההצגה השנייה של פונקציית גרין היא:
\[
G(x, x') = \begin{cases}
\dfrac{y_2(x') y_1(x)}{p_0(x') W(x')}, & x < x', \\[10pt]
\dfrac{y_1(x') y_2(x)}{p_0(x') W(x')}, & x > x'.
\end{cases}
\]

\subsection*{קבוע הרונסקיאן}

נוכיח שהמכנה \(p_0(x) W(x)\) הוא למעשה קבוע, ולא תלוי ב-\(x\) או \(x'\). זה יפשט את הביטוי. נגזור:
\[
\frac{d}{dx} \left[ p_0(x) W(x) \right] = p_0'(x) W(x) + p_0(x) W'(x).
\]
נחשב \(W'(x)\):
\[
W'(x) = \frac{d}{dx} \left[ y_1(x) y_2'(x) - y_1'(x) y_2(x) \right] = y_1'(x) y_2'(x) + y_1(x) y_2''(x) - y_1''(x) y_2(x) - y_1'(x) y_2'(x).
\]
האיברים הראשון והאחרון מתבטלים, ונשאר:
\[
W'(x) = y_1(x) y_2''(x) - y_1''(x) y_2(x).
\]
לכן:
\[
\frac{d}{dx} \left[ p_0(x) W(x) \right] = p_0' W + p_0 \left( y_1 y_2'' - y_1'' y_2 \right).
\]
כעת נזכור ש-\(y_1\) ו-\(y_2\) מקיימות את המשוואה ההומוגנית:
\[
\frac{d}{dx} \left( p_0 y_i' \right) + q y_i - \Lambda y_i = 0,
\]
כלומר
\[
p_0 y_i'' + p_0' y_i' + q y_i - \Lambda y_i = 0.
\]
נכפיל את המשוואה עבור \(y_2\) ב-\(y_1\), ואת המשוואה עבור \(y_1\) ב-\(y_2\), ונחסר:
\[
p_0 (y_1 y_2'' - y_1'' y_2) + p_0' (y_1 y_2' - y_1' y_2) = 0.
\]
אך \(y_1 y_2' - y_1' y_2 = W\), ולכן:
\[
p_0 (y_1 y_2'' - y_1'' y_2) + p_0' W = 0,
\]
כלומר
\[
\frac{d}{dx} \left[ p_0(x) W(x) \right] = 0.
\]
לכן \(p_0(x) W(x)\) קבוע. נהוג לבחור נקודה נוחה, למשל \(x = a\), ולסמן:
\[
p_0(x) W(x) = p_0(a) W(a) \equiv \text{const}.
\]

\subsection*{הצגה סופית}

נוכל כעת לכתוב את פונקציית גרין בצורה סופית ופשוטה יותר:
\[
G(x, x') = \begin{cases}
\dfrac{y_2(x') y_1(x)}{p_0(a) W(a)}, & x < x', \\[10pt]
\dfrac{y_1(x') y_2(x)}{p_0(a) W(a)}, & x > x'.
\end{cases}
\]

הביטוי הזה אינו תלוי בטור אינסופי, והוא מתאים לחישוב ישיר במקרים בהם אנחנו יכולים לפתור את המשוואה ההומוגנית. זוהי ההצגה השנייה והשימושית ביותר של פונקציית גרין.

\subsection*{סיכום ומשמעות פיזיקלית}

פונקציית גרין \(G(x, x')\) היא התגובה של המערכת לכוח נקודתי (דלתא) הממוקם ב-\(x = x'\). האינטגרל
\[
u(x) = \int_a^b G(x, x') f(x') \, dx'
\]
מייצג \textbf{סופרפוזיציה לינארית} של תגובות לכוחות נקודתיים: כל נקודה \(x'\) תורמת \(f(x')\) כפול התגובה \(G(x, x')\), ואנחנו מסכמים (מאינטגרלים) על כל הנקודות. זוהי המשמעות הפיזיקלית של פונקציית גרין.

בשיעור הבא נפתור דוגמה מפורטת ונראה כיצד להשתמש בשתי ההצגות של פונקציית גרין לפתרון בעיות קונקרטיות.
